
\subsection{Carry unsold stock through the assessed annual book}
The next composition supplies the annual execution step left open above. It
keeps the full tax-adjusted company book, including income tax payable,
assessment records, liability vintages, losses and earlier periods. The
standalone inventory method accepts 25 accounts, whereas this financial book
has 26. Removing tax payable or moving it into equity to reuse that smaller
schema would change the company. The annual extension instead validates and
updates the complete registered state.

The opening for this example comes from the checked retained-stock history,
with its planned whole-lot sale omitted. At 1 January 2027 it has two unsold
recovered units with carrying amount and tax base of 20, cash of 452 and tax
payable of $41/4$. New-period income starts at zero. Selling one unit for 15
on 2 January gives cash 467, inventory 70, a remaining unit with cost 10 and
recorded profit 5. Operating cash is 15. The payable remains $41/4$: an
inventory sale does not pay or reassess a prior tax liability.

The sale retains the complete assessment and liability histories. If it occurs
after a current assessment, that assessment's amounts remain recorded but its
status records subsequent activity. Any prior deferred measurement becomes
stale because it was bound to the earlier financial state. A sale inside the
already assessed interval is refused under the existing timing rule. New
activity is not silently inserted into a historical assessment.

Consider a hypothetical current assessment of 2 for the period containing the
first sale. It adds a separate liability, so total payable becomes $49/4$ and
recorded profit becomes 3. With explicitly supplied fresh deferred measurement
and the remaining carry conditions, a provisional boundary transfers that
profit to retained earnings and opens the next period with the unsold unit,
its cost of 10, cash 467 and the same liability vintages. Selling that unit for
15 on 2 January 2028 then gives cash 482 and recorded profit 5. That period's
operating cash is 15; the earlier receipt remains in the archived period.
The new-period payable is still $49/4$ until a payment or assessment changes it.

A separate sequence checks why fresh assessment matters. After a midyear
assessment and deferred measurement, an inventory sale records further
activity and invalidates the earlier measurement. Carry is unresolved at that
point. A fresh cumulative current assessment followed by a measurement bound
to the changed state permits the selected provisional carry; both earlier and
replacement assessments remain available. Reusing the old measurement after
that reassessment is invalid because its state digest no longer matches.
Attempting measurement before the required current reassessment instead stops
at that missing prerequisite. The two refusals answer different questions.

Ordinary operations continue through the protected annual evaluator. Their
financial states, reports and traces are unchanged; the new program profile
and request identities identify the extension. A full-quantity sale also
matches the old annual whole-lot sale on the complete resulting book, journal,
flows and report. Altering tax payable, an archived quantity or an allocation
fails complete-result verification even if the alteration is balanced.

The optional operation now has a checked annual execution meaning. Its matched
source declarations and annual-family dispatch still require separate
extensions. The current authoring and family profiles do not acquire support
merely because the financial transition exists. The examples remain conditional
on the stated inventory and tax assumptions; actual source and accounting
admission are unresolved.
